% Generated draft: requires mathematical and source review.
\documentclass[11pt]{article}
\usepackage[margin=25mm]{geometry}
\usepackage{fontspec}
\setmainfont{Latin Modern Roman}
\usepackage{amsmath,amssymb,mathtools,unicode-math}
\setmathfont{Latin Modern Math}
\usepackage{enumitem,needspace,xurl,hyperref}
\setlength{\parindent}{0pt}
\setlength{\parskip}{5pt}
\setlength{\emergencystretch}{3em}
\newcommand{\N}{\mathbb{N}}
\newcommand{\Z}{\mathbb{Z}}
\newcommand{\Q}{\mathbb{Q}}
\newcommand{\R}{\mathbb{R}}
\newcommand{\C}{\mathbb{C}}
\newcommand{\F}{\mathbb{F}}
\newcommand{\A}{\mathbb{A}}
\newcommand{\PP}{\mathbb{P}}
\newcommand{\E}{\mathbb{E}}
\newcommand{\eps}{\varepsilon}
\newcommand{\dd}{\,\mathrm{d}}
\begin{document}
\section*{UnsolvedMath ID: 30006115}
\textbf{Source problem code:} catalogue entry
\subsection{Definitions and mathematical statement}
\textbf{Formulation status.} The supplied record identifies an open-ended survey entitled \emph{Survey of Smooth Modular Representations}, rather than a recoverable standalone conjecture. Its statement status is \texttt{unrecoverable}. Consequently, no precise open mathematical assertion, with complete hypotheses and quantifiers, can be extracted from this record alone.

\textbf{Recoverable notation.} Whenever a group \(G\), a coefficient ring \(R\), a subgroup \(K\), and a category \(\operatorname{Mod}_R(G)\) of smooth representations have been specified, write
\[
V^K=\{v\in V:kv=v\text{ for every }k\in K\}.
\]
The condition that \(V\) is generated by its \(K\)-invariant vectors means
\[
V=\operatorname{span}_R\{gv:g\in G,\ v\in V^K\}.
\]
Let \(\mathcal C_K\) denote the full subcategory of such representations, and let \(\operatorname{Irr}_R(G)\) denote the isomorphism classes of irreducible objects of \(\operatorname{Mod}_R(G)\). These definitions explain the extracted notation; they do not determine the missing hypotheses on \(G\), \(R\), or \(K\).

\textbf{Extracted assertions requiring source recovery.} The fragment denotes an induced representation by \(P=\operatorname{ind}_K^G(1_R)\), with \(1_R\) apparently the trivial \(K\)-representation on \(R\). It reports that, if \(P\) is projective, the invariants assignment \(V\mapsto V^K\) gives a bijection of the schematic form
\[
\{[V]\in\operatorname{Irr}_R(G):V^K\ne0\}
\longleftrightarrow
\{[M]:M\text{ is a simple }H_R(G,K)\text{-module}\},
\]
and an equivalence
\[
\mathcal C_K\simeq\operatorname{Mod}(H_R(G,K)).
\]
The fragment also mentions stability under subrepresentations, an extension from projectivity to ``almost projectivity,'' and an Iwahori subgroup as an example. These are reported source claims, not assertions validated by this draft: the induction convention, coefficient hypotheses, definition of almost projectivity, Hecke algebra definition, module handedness, and precise role of stability under subrepresentations are absent or unclear. The opening claim about a convolution algebra being isomorphic to an opposite algebra has lost an essential referent.

A detached footnote appears to qualify some assertion as true for mod \(p\) representations of \(\operatorname{GL}_2(\mathbb Q_p)\) but false in general. Both the attachment of the footnote and the restoration of its damaged group notation require verification; no assertion above is assigned that qualification without evidence.

\textbf{Unresolved survey target.} The recoverable question is only what is known and unknown about smooth modular representations and their Hecke algebras. The supplied material does not identify a particular classification problem, universal assertion, or conjectural equivalence to formalize as an open problem.
\subsection{Short English statement}
This record contains a survey heading and a damaged passage about invariants and Hecke algebras. It does not supply a standalone mathematical problem. The displayed correspondences describe the extracted claims schematically; missing hypotheses, an undefined almost-projectivity condition, and a detached qualification prevent a precise theorem or conjecture statement. The supplied open label requires independent review.
\subsection{Sources}
\begin{itemize}
\item Local registry: \texttt{data/unsolvedmath/problems.json}, with \texttt{catalog/problems.json} as the supplementary registry; supplied complete record for permanent UnsolvedMath ID \(30006115\), exact source code \texttt{catalogue entry}, title \emph{Survey of Smooth Modular Representations}. No individual proposer is supplied; the proposed year is \(2024\).
\item Supplied report citation: \emph{Representations of \(p\)-adic Groups} (2025), Oberwolfach Reports, DOI 10.4171/owr/2024/54, pp. 3103--3178. \url{https://doi.org/10.4171/owr/2024/54}. This report citation is supplied provenance; the relevant passage has not been independently inspected for this draft.
\item Supplied literature source, labeled ``Vigneras survey report'': \url{https://publications.mfo.de/bitstream/handle/mfo/4270/OWR_2024_54.pdf?sequence=1}. The record attributes the material to an Oberwolfach workshop survey and its open-problem discussions; it supplies no more precise passage location or individual problem attribution.
\item Literature assessment date: \texttt{2026-08-22}. This is supplied dated curation, not independently verified. The supplied assessment labels the record open while explaining that the extracted text is an open-ended survey prompt rather than a resolved mathematical assertion. The supplied statement-verification note likewise identifies ``what is known and what is not known?'' as a broad survey heading, not a standalone problem. The record requests independent verification of time-sensitive status and cited claims.
\end{itemize}
\end{document}
